\scanpage{60}
% @ lines 1–3 / text
bornant cependant à \deleted{Spec} $S'_0=\deleted{\unsure{\mathrm{Sp}}}S-\{2\}$ dans les trois cas $(*)$ où l'opération \inserted{sur $S_0$} est fidèle, mais l'hom de grpes envisagé n'est pas épi.
% @ lines 4–12 / text
Comme les 11 catégories \inserted{fibrées} envisagées sont des gerbes (\unsure{fpqc|fppf} disons, ou ét), il s'ensuit que chacune d'elles est can. équiv. à la cat. des torseurs \inserted{à droite} sous le $\underline{\operatorname{Aut}}(\unsure{\mathscr X|\mathcal X|\mathfrak X}_0)$ correspondant, donc aussi à la cat. des torseurs à dr. sous $G_0$ (avec le grain de sel $*$). Ceci établira l'\uline{existence} d'un multidictionnaire, mais ne \inserted{nous} dispensera pas de la tâche de le développer en détail.

% @ line 13, heading / text
I Représentants typiques des onze cat. fibrées.
% @ lines 14–22, three columns / diagram
\diagram{$\begin{array}{@{}lll@{}}
% item line 14: right torsor
 1)\ C_{0_{S_0}}: & \multicolumn{2}{l}{\deleted{\unsure{R_0=|P_0=}}\text{Le torseur trivial }\inserted{R_0=(G_{S_0})_{\rm d}}\text{ sous }\mathrm{GP}(1)_{S_0}=G_{S_0}}\\[3pt]
% item lines 15–17: trivial group forms
 2)\ C_{1_{S_0}}: & \text{La forme triviale} &G_0=G_{S_0}\inserted{=\mathrm{GP}(1)}\text{ de }G_{S_0}\\
 \deleted{3)}\ C'_{1_{S_0}} &\rule[.4ex]{3.3cm}{.4pt} & \widetilde G_0=\unsure{\mathrm{Sl}|\mathrm{Gl}}(2)\text{. de }\mathrm{Sl}(2)\\
 \deleted{4)\ C^{\unsure{\prime|\prime\prime}}_{1_{S_0}}}&\deleted{\rule[.4ex]{3.3cm}{.4pt}}&\deleted{\unsure{\bar G_0}=\mathrm{Gl}(2)\text{ de }\mathrm{Gl}(2)}\\[4pt]
% item lines 18–20: trivial Lie forms
 \unsure{3|5})\ \left\{\begin{array}{l}C_{2_{S_0}}\\ C'_{2_{S_0}}\\ \deleted{C''_{\unsure{2|1}_{S_0}}}\end{array}\right.
 &\begin{array}{l}\rule[.4ex]{3.3cm}{.4pt}\\ \rule[.4ex]{3.3cm}{.4pt}\\\deleted{\rule[.4ex]{3.3cm}{.4pt}}\end{array}
 &\begin{array}{l}\unsure{\mathfrak g|g}_0=\mathfrak{gp}(1)\text{ de }\mathfrak{gp}(1)\\
 \widetilde{\mathfrak g}_0=\mathfrak{sl}(2)\text{ de }\mathfrak{sl}(2)\\
 \deleted{\unsure{\mathfrak L_0|\mathfrak g''_0}=\mathfrak{gl}(2)\text{ de }\mathfrak{gl}(2)}\end{array}\\[4pt]
% item line 21: matrix algebra
 4)\ C_{3_{S_0}}: &\multicolumn{2}{l}{\text{l'algèbre d'Azumaya triviale }M_0=\mathbb M_2(\mathcal O_{S_0})}\\
% item line 22: struck projective representative and beginning of following row
 \deleted{5)\left\{\begin{array}{l}C_{4_{S_0}}:\\ C'_{4_{S_0}}:\end{array}\right.}
 &\multicolumn{2}{l}{\deleted{X_0=\mathbb P^1_{S_0}=\mathbb P^1(\mathcal O_{S_0}^2)=\check{\mathbb P}(\mathcal O_{S_0}^2)}} 
\end{array}$}
% archive sheet number: 59
